\documentclass[reqno]{amsart}
\usepackage[T1]{fontenc}
\usepackage{amsmath,amssymb,amsthm}
\usepackage{xr-hyper}
\usepackage[hidelinks]{hyperref}
\externaldocument[ib-]{20261125T001100--inert-primes-and-estimate-b__workbook}
\externaldocument[cm-]{20261124T223100--complex-multiplication-and-inert-reduction__workbook}
\externaldocument[st-]{20261124T214500--supersingular-theta-divisibility__workbook}
\externaldocument[ip-]{20261124T214600--the-isotropic-plane__workbook}
\externaldocument[ht-]{20261125T013500--hecke-nondegeneracy-on-the-tree__workbook}
\externaldocument[as-]{20261125T021500--assembly-at-one-scale__workbook}

% Chapter number; \ref{chap} gives it.
\newcounter{chap}
% Chapter numbers, generated from 20261125T012000--grh-workbook__guide.org
% by scripts/make-chapters.py.  Do not edit: change the order of the
% index and rerun the script (./preview-site.sh and CI do this).
%
% Usage in a sheet, after \newcounter{chap}:
%   % Chapter numbers, generated from 20261125T012000--grh-workbook__guide.org
% by scripts/make-chapters.py.  Do not edit: change the order of the
% index and rerun the script (./preview-site.sh and CI do this).
%
% Usage in a sheet, after \newcounter{chap}:
%   % Chapter numbers, generated from 20261125T012000--grh-workbook__guide.org
% by scripts/make-chapters.py.  Do not edit: change the order of the
% index and rerun the script (./preview-site.sh and CI do this).
%
% Usage in a sheet, after \newcounter{chap}:
%   \input{chapters}\chapnumber{<sheet ID>}
\makeatletter
\providecommand\chapters@set[2]{\@namedef{chapters@#1}{#2}}
\chapters@set{20261125T001000}{1}% zeros-to-ramp-sums
\chapters@set{20261125T001100}{2}% inert-primes-and-estimate-b
\chapters@set{20261125T005000}{3}% galois-fourier-and-averaging
\chapters@set{20261125T005500}{4}% columns-from-a-sphere
\chapters@set{20261125T000000}{5}% slope-method
\chapters@set{20261126T030100}{6}% a-twist-for-each-prime
\chapters@set{20261125T003000}{7}% the-chain-and-the-square-root
\chapters@set{20261125T010000}{8}% rank-lemma-and-zero-estimate
\chapters@set{20261125T010500}{9}% poincare-series-and-near-orthogonality
\chapters@set{20261124T223000}{10}% elliptic-curves-for-the-proof
\chapters@set{20261125T009500}{11}% fock-model-and-theta-intertwiners
\chapters@set{20261124T223100}{12}% complex-multiplication-and-inert-reduction
\chapters@set{20261125T008000}{13}% one-cm-determinant
\chapters@set{20261124T214500}{14}% supersingular-theta-divisibility
\chapters@set{20261124T214600}{15}% the-isotropic-plane
\chapters@set{20261125T013700}{16}% pairings-on-a-formal-group-of-height-two
\chapters@set{20261125T013500}{17}% hecke-nondegeneracy-on-the-tree
\chapters@set{20261126T030200}{18}% the-bound-is-attained
\chapters@set{20261125T009000}{19}% kronecker-limit-and-chowla-selberg
\chapters@set{20261125T002000}{20}% definite-theta-lifts
\chapters@set{20261125T015500}{21}% siegel-weil-in-the-smallest-case
\chapters@set{20261125T002100}{22}% definite-theta-lifts-and-norms
\chapters@set{20261125T016500}{23}% hecke-operators-and-theta-lifts
\chapters@set{20261125T016700}{24}% divisibility-along-the-chain
\chapters@set{20261125T004000}{25}% denominators-and-sha
\chapters@set{20261125T004100}{26}% tate-twists-and-sha
\chapters@set{20261126T020430}{27}% the-minus-part-by-gauss-sums
\chapters@set{20261125T014100}{28}% congruence-modules
\chapters@set{20261125T014500}{29}% eisenstein-congruences-and-selmer-classes
\chapters@set{20261125T014700}{30}% pseudocharacters-and-the-length-bound
\chapters@set{20261125T011000}{31}% galois-representations-black-box
\chapters@set{20261126T030300}{32}% the-data
\chapters@set{20261126T020210}{33}% bessel-poincare-series-and-their-gram-operators
\chapters@set{20261126T030400}{34}% the-seesaw-identity
\chapters@set{20261126T030500}{35}% divisibility-of-the-composite-lift
\chapters@set{20261126T030600}{36}% theorem-a-short-forms
\chapters@set{20261126T030700}{37}% the-determinant-of-a-chain
\chapters@set{20261126T030800}{38}% the-cost-of-one-lift
\chapters@set{20261126T030900}{39}% denominators-of-projections-to-lifts
\chapters@set{20261126T031000}{40}% twists-with-few-congruences
\chapters@set{20261126T031100}{41}% theorem-c-and-the-main-inequality
\providecommand\chapnumber[1]{%
  \@ifundefined{chapters@#1}%
    {\PackageError{chapters}{Sheet #1 is not listed in the index}%
      {List it in 20261125T012000--grh-workbook__guide.org\MessageBreak
       and run scripts/make-chapters.py.}}%
    {\edef\thechap{\@nameuse{chapters@#1}}}}
\makeatother
\chapnumber{<sheet ID>}
\makeatletter
\providecommand\chapters@set[2]{\@namedef{chapters@#1}{#2}}
\chapters@set{20261125T001000}{1}% zeros-to-ramp-sums
\chapters@set{20261125T001100}{2}% inert-primes-and-estimate-b
\chapters@set{20261125T005000}{3}% galois-fourier-and-averaging
\chapters@set{20261125T005500}{4}% columns-from-a-sphere
\chapters@set{20261125T000000}{5}% slope-method
\chapters@set{20261126T030100}{6}% a-twist-for-each-prime
\chapters@set{20261125T003000}{7}% the-chain-and-the-square-root
\chapters@set{20261125T010000}{8}% rank-lemma-and-zero-estimate
\chapters@set{20261125T010500}{9}% poincare-series-and-near-orthogonality
\chapters@set{20261124T223000}{10}% elliptic-curves-for-the-proof
\chapters@set{20261125T009500}{11}% fock-model-and-theta-intertwiners
\chapters@set{20261124T223100}{12}% complex-multiplication-and-inert-reduction
\chapters@set{20261125T008000}{13}% one-cm-determinant
\chapters@set{20261124T214500}{14}% supersingular-theta-divisibility
\chapters@set{20261124T214600}{15}% the-isotropic-plane
\chapters@set{20261125T013700}{16}% pairings-on-a-formal-group-of-height-two
\chapters@set{20261125T013500}{17}% hecke-nondegeneracy-on-the-tree
\chapters@set{20261126T030200}{18}% the-bound-is-attained
\chapters@set{20261125T009000}{19}% kronecker-limit-and-chowla-selberg
\chapters@set{20261125T002000}{20}% definite-theta-lifts
\chapters@set{20261125T015500}{21}% siegel-weil-in-the-smallest-case
\chapters@set{20261125T002100}{22}% definite-theta-lifts-and-norms
\chapters@set{20261125T016500}{23}% hecke-operators-and-theta-lifts
\chapters@set{20261125T016700}{24}% divisibility-along-the-chain
\chapters@set{20261125T004000}{25}% denominators-and-sha
\chapters@set{20261125T004100}{26}% tate-twists-and-sha
\chapters@set{20261126T020430}{27}% the-minus-part-by-gauss-sums
\chapters@set{20261125T014100}{28}% congruence-modules
\chapters@set{20261125T014500}{29}% eisenstein-congruences-and-selmer-classes
\chapters@set{20261125T014700}{30}% pseudocharacters-and-the-length-bound
\chapters@set{20261125T011000}{31}% galois-representations-black-box
\chapters@set{20261126T030300}{32}% the-data
\chapters@set{20261126T020210}{33}% bessel-poincare-series-and-their-gram-operators
\chapters@set{20261126T030400}{34}% the-seesaw-identity
\chapters@set{20261126T030500}{35}% divisibility-of-the-composite-lift
\chapters@set{20261126T030600}{36}% theorem-a-short-forms
\chapters@set{20261126T030700}{37}% the-determinant-of-a-chain
\chapters@set{20261126T030800}{38}% the-cost-of-one-lift
\chapters@set{20261126T030900}{39}% denominators-of-projections-to-lifts
\chapters@set{20261126T031000}{40}% twists-with-few-congruences
\chapters@set{20261126T031100}{41}% theorem-c-and-the-main-inequality
\providecommand\chapnumber[1]{%
  \@ifundefined{chapters@#1}%
    {\PackageError{chapters}{Sheet #1 is not listed in the index}%
      {List it in 20261125T012000--grh-workbook__guide.org\MessageBreak
       and run scripts/make-chapters.py.}}%
    {\edef\thechap{\@nameuse{chapters@#1}}}}
\makeatother
\chapnumber{<sheet ID>}
\makeatletter
\providecommand\chapters@set[2]{\@namedef{chapters@#1}{#2}}
\chapters@set{20261125T001000}{1}% zeros-to-ramp-sums
\chapters@set{20261125T001100}{2}% inert-primes-and-estimate-b
\chapters@set{20261125T005000}{3}% galois-fourier-and-averaging
\chapters@set{20261125T005500}{4}% columns-from-a-sphere
\chapters@set{20261125T000000}{5}% slope-method
\chapters@set{20261126T030100}{6}% a-twist-for-each-prime
\chapters@set{20261125T003000}{7}% the-chain-and-the-square-root
\chapters@set{20261125T010000}{8}% rank-lemma-and-zero-estimate
\chapters@set{20261125T010500}{9}% poincare-series-and-near-orthogonality
\chapters@set{20261124T223000}{10}% elliptic-curves-for-the-proof
\chapters@set{20261125T009500}{11}% fock-model-and-theta-intertwiners
\chapters@set{20261124T223100}{12}% complex-multiplication-and-inert-reduction
\chapters@set{20261125T008000}{13}% one-cm-determinant
\chapters@set{20261124T214500}{14}% supersingular-theta-divisibility
\chapters@set{20261124T214600}{15}% the-isotropic-plane
\chapters@set{20261125T013700}{16}% pairings-on-a-formal-group-of-height-two
\chapters@set{20261125T013500}{17}% hecke-nondegeneracy-on-the-tree
\chapters@set{20261126T030200}{18}% the-bound-is-attained
\chapters@set{20261125T009000}{19}% kronecker-limit-and-chowla-selberg
\chapters@set{20261125T002000}{20}% definite-theta-lifts
\chapters@set{20261125T015500}{21}% siegel-weil-in-the-smallest-case
\chapters@set{20261125T002100}{22}% definite-theta-lifts-and-norms
\chapters@set{20261125T016500}{23}% hecke-operators-and-theta-lifts
\chapters@set{20261125T016700}{24}% divisibility-along-the-chain
\chapters@set{20261125T004000}{25}% denominators-and-sha
\chapters@set{20261125T004100}{26}% tate-twists-and-sha
\chapters@set{20261126T020430}{27}% the-minus-part-by-gauss-sums
\chapters@set{20261125T014100}{28}% congruence-modules
\chapters@set{20261125T014500}{29}% eisenstein-congruences-and-selmer-classes
\chapters@set{20261125T014700}{30}% pseudocharacters-and-the-length-bound
\chapters@set{20261125T011000}{31}% galois-representations-black-box
\chapters@set{20261126T030300}{32}% the-data
\chapters@set{20261126T020210}{33}% bessel-poincare-series-and-their-gram-operators
\chapters@set{20261126T030400}{34}% the-seesaw-identity
\chapters@set{20261126T030500}{35}% divisibility-of-the-composite-lift
\chapters@set{20261126T030600}{36}% theorem-a-short-forms
\chapters@set{20261126T030700}{37}% the-determinant-of-a-chain
\chapters@set{20261126T030800}{38}% the-cost-of-one-lift
\chapters@set{20261126T030900}{39}% denominators-of-projections-to-lifts
\chapters@set{20261126T031000}{40}% twists-with-few-congruences
\chapters@set{20261126T031100}{41}% theorem-c-and-the-main-inequality
\providecommand\chapnumber[1]{%
  \@ifundefined{chapters@#1}%
    {\PackageError{chapters}{Sheet #1 is not listed in the index}%
      {List it in 20261125T012000--grh-workbook__guide.org\MessageBreak
       and run scripts/make-chapters.py.}}%
    {\edef\thechap{\@nameuse{chapters@#1}}}}
\makeatother
\chapnumber{20261125T013700}
\renewcommand{\thesection}{\thechap.\arabic{section}}

\theoremstyle{definition}
\newtheorem{problem}{Problem}[section]
\theoremstyle{plain}
\newtheorem{theorem}{Theorem}[section]

\newcommand{\Z}{\mathbb{Z}}
\newcommand{\Q}{\mathbb{Q}}
\newcommand{\Oc}{\mathcal{O}}
\newcommand{\Cc}{\mathcal{C}}
\newcommand{\lgg}{\operatorname{lg}}
\newcommand{\sw}{\mathrm{sw}}

\setlength{\emergencystretch}{2em}

\title{Pairings on a formal group of height two}
\author{}
\date{November 25, 2026}

\begin{document}
\maketitle
\refstepcounter{chap}\label{chap}

\paragraph{Prerequisites.}
Chapter~\ref{cm-chap}, \emph{Complex multiplication and reduction at an inert prime}: formal groups, their logarithms and the invariant derivation \(D\) (Section~\ref{cm-sec:formal}), and the statement of Lubin--Tate theory with the model logarithm (Section~\ref{cm-sec:reduction}, Problem~\ref{cm-pr:model-log}).  Chapter~\ref{st-chap}, \emph{The invariant derivation at a supersingular prime}: the function \(E_p\) and Theorem~\ref{st-thm:single}.  Chapter~\ref{ip-chap}, \emph{The isotropic plane}, Theorem~\ref{ip-thm:main} and Section~\ref{ip-sec:questions}; Chapter~\ref{ht-chap} is mentioned once, for orientation.  Chapter~\ref{ib-chap}, Problem~\ref{ib-pr:kronecker}.  The ring \(\Z_{p^2}\) of integers of the unramified quadratic extension of \(\Q_p\), its Teichm\"uller roots of unity, and its Frobenius \(\sigma\), which satisfies \(\sigma(a) \equiv a^p \pmod p\).  Legendre's formula \(v_p(n!) = (n - s_p(n))/(p-1)\), where \(s_p(n)\) is the sum of the base-\(p\) digits of \(n\).

\medskip
Chapter~\ref{ip-chap} found the gain at the odd levels only for \(p \equiv 3 \pmod 8\), because it needed a global element of the CM order acting by a non-unit on one eigenline.  The argument needs the gain at every inert prime.  What structure, visible at \(p\) alone, forces a theta kernel to be divisible by \(p^{\alpha p/(p^2-1)}\), that is by \(p^{\sum_{i \text{ odd}} \alpha/p^i}\), after \(\alpha\) derivatives?  In the argument a theta kernel at an inert prime has the form \(\exp(-\beta(x,y))\,U(x,y)\) with \(U\) integral, where \(x\) and \(y\) are parameters of a formal group of height two and \(\beta\) is a pairing, biadditive and semilinear for Frobenius up to series with integral exponential.  This chapter classifies such pairings and derives the divisibility from the normal form, in the simplest setting: one variable \(x\) and one partner variable \(y\).

\section{The formal module of height two}\label{sec:module}

Throughout, \(p\) is an odd prime, \(\Oc := \Z_{p^2}\), \(K := \Q_{p^2}\), \(\sigma\) is the Frobenius of \(K\), and \(v\) is the valuation with \(v(p) = 1\).  A power series over \(K\), or over a finite extension of \(K\), is \emph{integral} if all its coefficients have valuation \(\ge 0\).  Let
\[
  \lgg(X) := \sum_{j \ge 0} \frac{X^{p^{2j}}}{p^j},
  \qquad
  \lgg^*(X) := \frac1p \lgg(X^p) = \sum_{j \ge 0} \frac{X^{p^{2j+1}}}{p^{j+1}} .
\]
In the notation of Chapter~\ref{cm-chap}, \(\lgg = \ell_p\) is the model logarithm for the uniformizer \(p\) of \(\Oc\), with \(q = p^2\).

We use Hazewinkel's functional-equation lemma \cite[\S2.2]{Ha}, \cite[(2.5)--(2.8)]{Ha79} in the following case.  Let \(f(X) \in K[[X]]\) satisfy \(f(X) - p^{-1} f(X^{p^2}) \in \Oc[[X]]\) and \(f(X) \equiv X\) modulo degree \(2\).  Then: \(f^{-1}(f(X) + f(Y))\) has coefficients in \(\Oc\); if \(g(X) \in X K[[X]]\) also satisfies \(g(X) - p^{-1} g(X^{p^2}) \in \Oc[[X]]\), then \(f^{-1}(g(X)) \in \Oc[[X]]\); and for \(\alpha, \beta \in X\Oc[[X]]\) and \(r \ge 1\), \(\alpha \equiv \beta \pmod{p^r}\) if and only if \(f(\alpha) \equiv f(\beta) \pmod{p^r}\).

Let \(F(X,Y) := \lgg^{-1}(\lgg X + \lgg Y)\) and \([\alpha](X) := \lgg^{-1}(\alpha \lgg X)\) for \(\alpha \in \Oc\).

\begin{problem}\label{pr:LT}
\(F \in \Z_p[[X,Y]]\) and \([\alpha] \in \Oc[[X]]\), so \(F\) is a formal \(\Oc\)-module.  Modulo degree \(p^2 + 1\),
\[
  F(X,Y) \equiv X + Y - \frac1p \sum_{0 < k < p^2} \binom{p^2}{k} X^k Y^{p^2-k} .
\]
For \(\zeta \in \Oc\) with \(\zeta^{p^2-1} = 1\), \([\zeta](X) = \zeta X\) and \(F(\zeta X, \zeta Y) = \zeta F(X,Y)\).
\end{problem}

\begin{problem}\label{pr:height}
\([p](X) \equiv X^{p^2} \pmod p\).  So \(F\) has height two.
\end{problem}

\begin{problem}\label{pr:twist}
Let \(u\) be a root of unity of order \(2(p^2-1)\), so that \(u^{p^2-1} = -1\) and \(u \in \Z_{p^4}\).  Then \(\ell_{-p}(uX) = u \lgg(X)\).  Hence for \(p \equiv 3 \pmod 4\) the formal group of \(y^2 = x^3 - x\) is isomorphic to \(F\) over \(\Z_{p^4}\).
\end{problem}

For a series \(f\) over \(K\) let \(f^\sigma\) be the series obtained by applying \(\sigma\) to the coefficients, and \(f^\phi(X) := f^\sigma(X^p)\), where \(X^p\) raises every variable to the power \(p\).

\begin{problem}\label{pr:frob}
For \(f\) integral over \(\Oc\) in any number of variables, \(f^\phi \equiv f^p \pmod p\).  In particular \(F(X^p, Y^p) \equiv F(X,Y)^p \pmod p\).
\end{problem}

\section{Pairings and their normal form}\label{sec:pairings}

For series \(f\), \(g\) over \(K\) write \(f \simeq g\) if \(f - g\) has no constant term and \(\exp(f - g)\) is integral.  The basic tool is Dwork's lemma \cite[\S IV.2]{Ko}, in the following form \cite[Lemma~9.1(i)]{Adr}, which has the same proof: for a series \(\Psi\) over \(K\) in finitely many variables without constant term,
\[
  \exp\Psi \text{ is integral} \iff \Psi^\phi - p\Psi \in p\,\Oc[[X]] .
\]
For a series \(H\) in two variables let \(H^\sw(X,Y) := H(Y,X)\) and
\[
  (d_F H)(X_0, X_1, X_2) := H\bigl(F(X_0,X_1), X_2\bigr) - H(X_0, X_2) - H(X_1, X_2),
\]
and for a series \(f\) in one variable let \(\partial_F f := f(F(X_0,X_1)) - f(X_0) - f(X_1)\).  A \emph{pairing} is a series \(\beta(X,Y)\) over \(K\) such that:
(i) every monomial of \(\beta\) contains both variables;
(ii) \(d_F\beta \simeq 0\) and \(d_F\beta^\sw \simeq 0\);
(iii) \(\beta(\zeta X, Y) \simeq \beta(X, \zeta^p Y)\) for every \(\zeta \in \Oc\) with \(\zeta^{p^2-1} = 1\).
Condition (ii) says that \(\exp\beta\) is biadditive up to integral cocycles.  Condition (iii) is semilinearity for Frobenius, since \(\zeta^p = \sigma(\zeta)\).

\begin{problem}\label{pr:simeq}
The relation \(\simeq\) is an equivalence relation, stable under sums and under substitution of integral series without constant term for the variables.  If \(f \simeq g\), then \(d_F f \simeq d_F g\).
\end{problem}

\begin{problem}\label{pr:average}
(i) \(\beta(\zeta X, Y) = \beta(X, \zeta^p Y)\) for all \(\zeta\) as in (iii) if and only if every monomial \(X^m Y^n\) of \(\beta\) has \(m \equiv pn \pmod{p^2-1}\).
(ii) Let \(\beta\) be a pairing, \(N := p^2 - 1\), \(\zeta\) a primitive \(N\)th root of unity, and \(G := N^{-1}\sum_{k=0}^{N-1} \beta(\zeta^k X, \zeta^{-pk} Y)\).  Then \(G\) is a pairing, \(G \simeq \beta\), and \(G\) satisfies the condition of (i).
\end{problem}

For \(c, c' \in K\) let
\[
  A(X,Y) := \sum_{0 \le i \le j} \frac{X^{p^{2i}}\, Y^{p^{2j+1}}}{p^{\,i+j+1}},
  \qquad
  Q_{c,c'}(X,Y) := c\,A(X,Y) - c'\,A(Y,X) .
\]
In every monomial of \(A\) the exponent of \(Y\) is the exponent of \(X\) times an odd power of \(p\).

\begin{problem}\label{pr:Q-semilinear}
For all \(c, c'\), \(Q_{c,c'}(\zeta X, Y) = Q_{c,c'}(X, \zeta^p Y)\) for every \(\zeta\) as in (iii).  No monomial of \(\lgg(X)\lgg(Y)\) satisfies the congruence of Problem~\ref{pr:average}(i).
\end{problem}

\begin{problem}\label{pr:reindex}
\(A(X^p, Y^p) = p\bigl(\lgg^*(X)\lgg(Y) - A(Y,X)\bigr)\).  Hence
\[
  Q_{c',c}(X^p, Y^p) - p\,Q_{c,c'}(X,Y) = p\bigl(c'\lgg^*(X)\lgg(Y) - c\lgg(X)\lgg^*(Y)\bigr).
\]
\end{problem}

\begin{problem}\label{pr:unique}
If \(c \in K\) and \(Q_{c,\sigma c} \simeq 0\), then \(c = 0\).
\end{problem}

\begin{problem}\label{pr:c-residue}
If \(c \in \Oc\), \(\beta\) is a pairing and \(\beta \simeq Q_{c,\sigma c}\), then \(p \cdot [X Y^p]\,\beta \in \Oc\) and \(p \cdot [XY^p]\,\beta \equiv c \pmod p\).
\end{problem}

Is \(Q_{c,\sigma c}\) itself a pairing?  It is not integral, and it is not biadditive on the nose.  Problems~\ref{pr:taylor}--\ref{pr:Q-pairing} show that it is biadditive up to \(\simeq\).

\begin{problem}\label{pr:taylor}
Let \(\Phi\) be a series over \(K\) whose partial derivative \(\partial\Phi/\partial X_e\) is integral, let \(\varphi\) be a tuple of integral series without constant terms, and let \(V\) be an integral series without constant term.  Then \(\Phi(\varphi + pV\mathbf{e}_e) - \Phi(\varphi) \in p\,\Oc[[X]]\), where \(\mathbf{e}_e\) is the unit vector of the variable \(X_e\).
\end{problem}

\begin{problem}\label{pr:dlg}
\(\partial_F \lgg = 0\), and \(\partial_F \lgg^*\) is integral.
\end{problem}

\begin{problem}\label{pr:dX}
For \(c, c' \in \Oc\), the series \(\partial_X\bigl(Q_{c,c'}(X,Y) - c\lgg(X)\lgg^*(Y)\bigr)\) is integral.
\end{problem}

\begin{problem}\label{pr:Q-pairing}
For \(c \in \Oc\), \(Q_{c,\sigma c}\) is a pairing.
\end{problem}

Conversely, every pairing has this form.

\begin{theorem}[{\cite[Thm.~9.3]{Adr}}]\label{thm:normal}
For every pairing \(\beta\) there is a unique \(c \in \Oc\) with \(\beta \simeq Q_{c,\sigma c}\).
\end{theorem}

Problems~\ref{pr:unique} and~\ref{pr:Q-pairing} give the uniqueness and show that every \(Q_{c,\sigma c}\) occurs.  For the existence, Adriazola replaces \(\beta\) by the average \(G\) of Problem~\ref{pr:average}, finds \(c\) from the coefficient of \(X\) in \(G\), as in Problem~\ref{pr:c-residue}, and applies Dwork's lemma to \(G - Q_{c,\sigma c}\), using Problem~\ref{pr:reindex}; we do not reproduce the argument.  The exactly bilinear pairing \(\lgg(X)\lgg(Y)\) is excluded by semilinearity (Problem~\ref{pr:Q-semilinear}).  So the denominators of a pairing sit only on the monomials \(X^{p^{2i}} Y^{p^{2j+1}}\) and \(Y^{p^{2i}} X^{p^{2j+1}}\), at odd shifts.

\section{Odd shifts and the credit}\label{sec:credit}

How much do these denominators cost after \(\alpha\) derivatives?  Let \(\varpi_p\) be a root of \(\varpi_p^{p^2-1} = p\), and let \(\Oc' := \Oc[\varpi_p]\), so that \(v(\varpi_p) = \rho\) with
\[
  \rho := \frac{1}{p^2-1},
  \qquad
  \gamma_p := \frac{p}{p^2-1} .
\]

\begin{problem}\label{pr:gamma}
\(\gamma_p = \sum_{e \ge 0} p^{-(2e+1)}\) and \(\gamma_p + \rho = 1/(p-1)\).  If \(n = \sum_k d_k p^k\) in base \(p\), then
\[
  \gamma_p n - E_p(n) = \gamma_p \sum_{k \text{ even}} d_k + \rho \sum_{k \text{ odd}} d_k,
  \qquad
  v(n!) - \rho n = \gamma_p n - \frac{s_p(n)}{p-1} .
\]
Hence \(v(n!) - \rho n \le E_p(n) \le \gamma_p n\).
\end{problem}

So a derivative of order \(n\) in a formal parameter gains \(v(n!)\), and after the rescaling \(x \mapsto \varpi_p x\) the gain is \(\gamma_p n\) up to the digits of \(n\); for the formal group of Chapter~\ref{st-chap}, a derivative of order \(n\) in the logarithm gains \(E_p(n)\) on integral series (Theorem~\ref{st-thm:single}).  Both are the credit \(\gamma_p n = \sum_{i \text{ odd}} n/p^i\) up to digits.

\begin{problem}\label{pr:rescaled}
In \(Q_{c,c'}(\varpi_p X, \varpi_p Y)\) the coefficient of \(X^{p^{2i}} Y^{p^{2j+1}}\) (\(i \le j\)) has valuation \(v(c) + (p^{2i} + p^{2j+1})\rho - (i+j+1)\).  For \(i = j = 0\) this is \(v(c) + \frac1{p-1} - 1\); otherwise it is at least \(v(c) + \frac{1}{p-1} + p - 2\).
\end{problem}

\begin{problem}\label{pr:exp-integral}
If every coefficient of a series \(\Psi\) without constant term has valuation \(\ge 1/(p-1)\), then \(\exp\Psi\) is integral over \(\Oc'\).  For every series \(h\) over \(\Oc'\) without constant term, \(\exp(\varpi_p^{p+1} h/p) = \sum_{k \ge 0} (h/p)^k\, g_k\) with \(g_k \in \Oc'\).
\end{problem}

For a series \(h\) without constant term, let \(\Cc(h)\) be the set of series \(\sum_{k \ge 0} p^{-k} h^k g_k\) with \(g_k\) integral over \(\Oc'\), the sum being finite in each degree.  It is stable under multiplication by integral series.  Theorem~\ref{thm:normal} and Problems~\ref{pr:rescaled}--\ref{pr:exp-integral} prove the following.

\begin{theorem}\label{thm:class}
Let \(\beta\) be a pairing with \(\beta \simeq Q_{c,\sigma c}\), let \(U \in \Oc[[x,y]]\), and let \(\Theta := \exp(-\beta(x,y))\, U(x,y)\).  Then
\[
  \Theta(\varpi_p x, \varpi_p y) \in \Cc(h),
  \qquad
  h := \sigma(c)\, x^p y - c\, x y^p .
\]
\end{theorem}

\section{Two local estimates}\label{sec:estimates}

Fix \(\Theta\), \(\beta\), \(c\), \(h\) as in Theorem~\ref{thm:class}.  We bound the mixed derivatives at \(0\) of \(\Theta\), of order \(a\) in \(x\) and \(w\) in \(y\), in two ways: in the formal parameters, with \(\partial_x^a \partial_y^w\), and in the logarithm coordinates, with \(D_x^a D_y^w\), where \(D_x := \lgg'(x)^{-1}\partial_x\) is the invariant derivation of Chapter~\ref{cm-chap}.  Thus \((D_x^a D_y^w \Theta)(0) = a!\, w!\, [z^a t^w]\, \Theta(\lgg^{-1} z, \lgg^{-1} t)\).  We call \(x\) the \emph{top} and \(y\) its \emph{partner}.

\begin{problem}\label{pr:estA}
If \(G(\varpi_p x, \varpi_p y) \in \Cc(h)\), then
\[
  v\bigl((\partial_x^a\partial_y^w G)(0)\bigr) \ \ge\ v(a!) + v(w!) - \rho(a+w) - \min\Bigl(w,\ \Bigl\lfloor \frac ap \Bigr\rfloor + \Bigl\lfloor \frac wp \Bigr\rfloor\Bigr).
\]
\end{problem}

In the formal parameters the multiplication by \(p\) on the partner, \(y \mapsto [p](y)\), is not linear.  In the logarithm coordinates it is \(t \mapsto pt\), and its effect on \(D_y^w\) is the factor \(p^w\).

\begin{problem}\label{pr:hecke}
\(\beta(x, [p]y) \simeq p\,\beta(x,y)\), and \(\exp\bigl(-p\,Q_{c,\sigma c}(\varpi_p x, \varpi_p y)\bigr)\) is integral over \(\Oc'\).
\end{problem}

\begin{problem}\label{pr:logcoord}
If \(G(\varpi_p x, \varpi_p y)\) is integral over \(\Oc'\), then \(v\bigl((D_x^a D_y^w G)(0)\bigr) \ge v(a!) + v(w!) - \rho(a+w)\).
\end{problem}

Applying Problem~\ref{pr:estA} to \(\Theta\), and Problem~\ref{pr:logcoord} to \(\Theta(x, [p]y)\), which is allowed by Problem~\ref{pr:hecke}, gives the following.

\begin{theorem}\label{thm:two}
For all \(a, w \ge 0\):
\begin{align*}
  \text{(A)}\quad & v\bigl((\partial_x^a\partial_y^w \Theta)(0)\bigr) \ \ge\ v(a!) + v(w!) - \rho(a+w) - \min\bigl(w,\ \lfloor a/p \rfloor + \lfloor w/p \rfloor\bigr), \\
  \text{(B)}\quad & v\bigl(p^w (D_x^a D_y^w \Theta)(0)\bigr) \ \ge\ v(a!) + v(w!) - \rho(a+w) .
\end{align*}
\end{theorem}

By Problem~\ref{pr:gamma}, both right sides are at least \(\gamma_p a - w - (s_p(a) + s_p(w))/(p-1)\): the top earns the credit \(\gamma_p a = \sum_{i \text{ odd}} a/p^i\), and the partner pays at most one factor of \(p\) per derivative.  In (B) the payment is the explicit factor \(p^w\), the effect of \([p]\) on the partner; in the argument this factor becomes a Hecke operator, whose elementary divisors are the loss computed in Chapter~\ref{ht-chap}.  The bounds use only the formal \(\Oc\)-module \(F\) and the semilinearity (iii) for the Teichm\"uller roots of unity, as in hypothesis (iii) of \cite[Thm.~9.3]{Adr}, not any global element of a CM order.

How are the \(w\) factors of \(p\) spent?  For \(\Theta = \exp(-Q_{c,\sigma c})\) and \(w < p\) they are spent once at each odd level.

\begin{problem}\label{pr:mixed}
Let \(c \in \Oc\), \(1 \le w < p\), \(e \ge 0\), and \(a := p^{2e+1} w\).  Then
\[
  \bigl(\partial_x^a \partial_y^w \exp(-Q_{c,\sigma c}(x,y))\bigr)(0) = \frac{a!\, \sigma(c)^w}{p^{(e+1)w}} ,
\]
of valuation \(v(a!) - (e+1)w\) when \(c\) is a unit.  (For \(p = 3\), \(w = 1\), \(e = 1\): \(27!\,\sigma(c)/9\), of valuation \(11\).)
\end{problem}

So the term \(x^{p^{2e+1}} y / p^{e+1}\) of \(A(y,x)\) charges \(w\) at each of the levels \(1, 3, \dotsc, 2e+1\); this is the charge of \(w\) at each odd level in the bound \(\sum_{i \text{ odd}} (\lfloor a/p^i \rfloor - w)_+\) of Theorem~\ref{ip-thm:main}.  The odd powers \(p^i\) are exactly those at which \(1 - \chi_d(p^i) = 2\) for \(p\) inert (Problem~\ref{ib-pr:kronecker}(ii)), and the credit \(\gamma_p\alpha\) at a top of weight \(\alpha\) is matched against them in \cite[\S3.2]{Adr}.

In the argument, \(x\) and \(y\) are the parameters along a hyperbolic pair \(f, f'\) of the hermitian form at \(p\), as in Section~\ref{ip-sec:questions}, transported to \(F\) by the Lubin--Tate isomorphism \cite[Prop.~8.11]{Adr}, as in Problem~\ref{pr:twist}.  For the CM curve used there and every prime \(p \ge 5\) inert in its CM field, Adriazola shows that the coboundary of the logarithm of the reduced theta function is a pairing \cite[Thm.~8.12]{Adr}, that the theta kernels are \(c_0 \exp(-\sum \beta(x,y))\, U\) with \(U\) integral \cite[Thm.~8.21]{Adr}, and, with many tops and partners and with minors in place of single derivatives, proves lower bounds of the kind of Theorem~\ref{thm:two} \cite[Thm.~9.6]{Adr}; its part~(iii), with an explicit power of \(p\) attached to the partners, corresponds to~(B).  That the constant \(c_0\) is integral is assumed there: in \cite{Adr} this is \cite[Hyp.~11.42, (N3)]{Adr}; see the ledger in Section~\ref{as-sec:map} of Chapter~\ref{as-chap}.  The Lean formalization \cite{Lean} does not assume it: it proves an integral kernel law of this shape for the kernels it uses, by a different route, in which the exponent is a torus average of a coboundary of the logarithm of the sigma function of one CM curve, and the integrality comes from a factorization of the fourth tensor power of the kernel into sigma functions.

\section*{Hints}

\subsection*{\S\ref{sec:module}}
\begin{description}
\item[\ref{pr:LT}] Apply the functional-equation lemma with \(f = \lgg\) and \(g = \alpha\lgg\).  Modulo degree \(p^2+1\), \(\lgg(X) \equiv X + X^{p^2}/p\).  For \(\zeta\), use \(\zeta^{p^{2j}} = \zeta\).
\item[\ref{pr:height}] The last part of the functional-equation lemma, with \(\lgg([p]X) - \lgg(X^{p^2}) = pX\).
\item[\ref{pr:twist}] \(u^{p^{2j}} = (-1)^j u\).  Then use Problem~\ref{cm-pr:model-log}.
\item[\ref{pr:frob}] \((f + g)^p \equiv f^p + g^p\) and \(\sigma(a) \equiv a^p\).
\end{description}

\subsection*{\S\ref{sec:pairings}}
\begin{description}
\item[\ref{pr:simeq}] An integral series with constant term \(1\) has an integral inverse.  For the last statement, express \(\exp(d_F(f - g))\) through \(\exp(f - g)\).
\item[\ref{pr:average}] (i) Compare \(\zeta^m\) with \(\zeta^{pn}\).  (ii) Apply (iii) with \(\zeta^k\) in place of \(\zeta\), and use the binomial series \((1+T)^{1/N}\).
\item[\ref{pr:Q-semilinear}] \(\zeta^{p^2} = \zeta\).
\item[\ref{pr:reindex}] Expand \(\lgg^*(X)\lgg(Y)\) as a double sum and split it according to the relative size of the two indices.
\item[\ref{pr:unique}] Dwork's lemma and Problem~\ref{pr:reindex}.  Look at the coefficients of \(X^p Y^{p^{2j}}\).
\item[\ref{pr:c-residue}] By condition (i) of a pairing, the coefficient of \(XY^n\) in \(\exp\Psi\) equals that of \(\Psi\).
\item[\ref{pr:taylor}] Reduce to a monomial \(aX_e^k M\) with \(v(ka) \ge 0\).  Use \(\binom kn = \frac kn \binom{k-1}{n-1}\) and \(v(p^n/n) \ge 1\).
\item[\ref{pr:dlg}] For \(\lgg^*\), use Problems~\ref{pr:frob} and~\ref{pr:taylor}, and that \(\lgg'\) is integral.
\item[\ref{pr:dX}] Compare each exponent with its denominator.
\item[\ref{pr:Q-pairing}] The second variable follows from the first, since \(Q_{c,\sigma c}^\sw = Q_{c'',\sigma c''}\) with \(c'' = -\sigma c\).  For the first, apply Dwork's lemma to \(d_F Q_{c,\sigma c}\).  The series \((d_FQ)^\phi\) differs from \(d_F(Q^\phi)\) by evaluating \(Q^\sigma = Q_{\sigma c, c}\) at \(F(X_0^p, X_1^p)\) instead of \(F(X_0,X_1)^p\).  Split \(Q_{\sigma c,c}\) as in Problem~\ref{pr:dX}, and use Problems~\ref{pr:reindex}, \ref{pr:taylor}, \ref{pr:dlg} and \(\lgg^*(Y^p) = \lgg(Y) - Y\).
\end{description}

\subsection*{\S\ref{sec:credit}}
\begin{description}
\item[\ref{pr:gamma}] \(\gamma_p n - E_p(n) = \sum_{i \text{ odd}} \{n/p^i\}\), the fractional parts.  Legendre's formula.
\item[\ref{pr:exp-integral}] \(v(k!) < k/(p-1)\) for \(k \ge 1\).
\end{description}

\subsection*{\S\ref{sec:estimates}}
\begin{description}
\item[\ref{pr:estA}] \(\partial_x^a\partial_y^w\) commutes with the rescaling up to \(\varpi_p^{a+w}\).  Count the factors \(x^p y\) and \(x y^p\) in a monomial of \(h^k\) of degree \(a\) in \(x\) and \(w\) in \(y\).
\item[\ref{pr:hecke}] \([p]\) is the \(p\)-fold sum for \(F\).  Problems~\ref{pr:rescaled} and~\ref{pr:exp-integral}.
\item[\ref{pr:logcoord}] The series \(\varpi_p^{-1}\lgg(\varpi_p T)\) has coefficients in \(\Z_p\).
\item[\ref{pr:mixed}] Only the terms of \(A(y,x)\) with \(i = 0\) have degree less than \(p\) in \(y\).  Then use the uniqueness of base-\(p\) digits.
\end{description}

\begin{thebibliography}{Ha79}
\bibitem[Adr]{Adr} J.~Adriazola, \emph{The Alp\"oge proof of the Riemann hypothesis}, manuscript, revised 24 November 2026, \url{https://github.com/jimjamzola/RH-Alpoge-Reconstruction}.
\bibitem[Ha]{Ha} M.~Hazewinkel, \emph{Formal Groups and Applications}, Academic Press, 1978.
\bibitem[Ha79]{Ha79} M.~Hazewinkel, \emph{On formal groups.  The functional equation lemma and some of its applications}, Ast\'erisque \textbf{63} (1979), 73--82.
\bibitem[Ko]{Ko} N.~Koblitz, \emph{\(p\)-adic Numbers, \(p\)-adic Analysis, and Zeta-Functions}, 2nd ed., Grad.\ Texts in Math.\ 58, Springer, 1984.
\bibitem[Lean]{Lean} The Lean community, \emph{riemann-hypothesis}, a Lean formalization of Alp\"oge's argument, \url{https://github.com/leanprover-community/riemann-hypothesis}.
\end{thebibliography}

\end{document}
